\section{CPU Implementation}

\subsection{Reference Implementations}

For all four ciphers --- \textbf{SIMECK, SIMON, TWINE, and WARP} --- correct reference implementations were developed based on their standard specifications.

\begin{itemize}
    \item \textbf{SIMECK and SIMON}: Implemented as ARX-based Feistel structures using bitwise rotation, AND, and XOR operations.
    \item \textbf{TWINE}: Implemented as a lightweight Substitution-Permutation Network (SPN) operating on 4-bit nibbles.
    \item \textbf{WARP}: Implemented as a substitution-permutation cipher with round constants and linear diffusion.
\end{itemize}

Each implementation includes:
\begin{itemize}
    \item A complete \textbf{key schedule}
    \item Correct \textbf{round transformations}
    \item Deterministic plaintext initialization for reproducibility
\end{itemize}

\subsection{Single-Threaded Implementation}

A baseline single-threaded version was implemented for each cipher. These implementations process one block at a time using sequential loops over all rounds.

These serve as:
\begin{itemize}
    \item A \textbf{correctness reference}
    \item A \textbf{performance baseline}
\end{itemize}

Execution time was measured using high-resolution timers (\texttt{chrono} or \texttt{clock}), and throughput was computed as:

\begin{equation}
\text{Throughput} = \frac{\text{Total Data Processed}}{\text{Execution Time}}
\end{equation}

\subsection{Multi-Threaded CPU Implementation}

Parallel CPU implementations were developed using \textbf{OpenMP}.

\subsubsection{Parallelization Strategy}

Encryption was parallelized across plaintext blocks:

\begin{verbatim}
#pragma omp parallel for
for (int i = 0; i < N; i++)
    encrypt(...)
\end{verbatim}

\subsubsection{Bit-Slicing Optimization}

For \textbf{TWINE} and \textbf{WARP}, a bit-sliced implementation was used:
\begin{itemize}
    \item 32 plaintext blocks are processed simultaneously
    \item Each bit position is stored across a 32-bit word
    \item Enables SIMD-style parallelism on CPU
\end{itemize}

This approach improves computational efficiency and reduces memory overhead.

\subsection{CPU Configuration}

All CPU experiments were conducted on the following system:

\begin{itemize}
    \item \textbf{Processor:} AMD Ryzen 7 5800H
    \item \textbf{Cores/Threads:} 8 cores / 16 threads
    \item \textbf{RAM:} 16 GB
    \item \textbf{Architecture:} x86\_64
    \item \textbf{Operating System:} Windows 11
    \item \textbf{Compiler/IDE:} Visual Studio 2022 (MSVC v143)
\end{itemize}

\subsection{Metrics Collected}

For each cipher and configuration:
\begin{itemize}
    \item Execution time (seconds)
    \item Throughput (GB/s)
    \item Speedup over single-threaded baseline
\end{itemize}

\section{GPU Implementation (CUDA)}

\subsection{Parallelization Strategy}

CUDA-based implementations were developed for all ciphers.

\subsubsection{Thread Mapping}

\begin{itemize}
    \item Each thread processes:
    \begin{itemize}
        \item One block (SIMECK, SIMON)
        \item One batch of 32 blocks (TWINE, WARP using bit-slicing)
    \end{itemize}
\end{itemize}

\subsubsection{Kernel Design}

A grid-stride loop was used for scalability:

\begin{verbatim}
for (int i = idx; i < N; i += stride)
    encrypt(...)
\end{verbatim}

\subsubsection{Thread Configuration}

\begin{itemize}
    \item \textbf{Block size:} 256 threads
    \item \textbf{Grid size:} $\lceil N / 256 \rceil$
\end{itemize}

\textbf{Justification:}
\begin{itemize}
    \item 256 threads aligns with warp scheduling (multiples of 32)
    \item Ensures high occupancy on the GPU
    \item Grid-stride loops allow efficient scaling for large inputs
\end{itemize}

\subsection{Memory Management}

\subsubsection{Host-Device Transfers}

The following were explicitly measured:
\begin{itemize}
    \item Host-to-Device (H2D) transfer time
    \item Device-to-Host (D2H) transfer time
\end{itemize}

\subsubsection{Memory Placement}

\begin{itemize}
    \item \textbf{Global Memory:}
    \begin{itemize}
        \item Plaintext and ciphertext buffers
        \item Large data structures
    \end{itemize}

    \item \textbf{Constant Memory:}
    \begin{itemize}
        \item Round keys (SIMECK, SIMON)
        \item Round constants and permutation tables (WARP)
    \end{itemize}

    \item \textbf{Registers:}
    \begin{itemize}
        \item Temporary variables during encryption rounds
    \end{itemize}
\end{itemize}

\textbf{Justification:}
\begin{itemize}
    \item Constant memory provides fast broadcast access for read-only data
    \item Global memory is necessary for large datasets
    \item Bit-slicing increases arithmetic intensity, improving GPU utilization
\end{itemize}

\subsection{GPU Configuration}

All GPU experiments were conducted on:

\begin{itemize}
    \item \textbf{GPU:} NVIDIA RTX 3050 Ti (95W TDP)
    \item \textbf{CUDA Version:} CUDA 12.3
    \item \textbf{Compiler/IDE:} Visual Studio 2022 (MSVC v143)
\end{itemize}

\section{Correctness Verification}

Correctness was verified through the following methods:

\begin{itemize}
    \item \textbf{CPU vs GPU comparison:}
    \begin{itemize}
        \item Outputs (ciphertexts) were matched for identical inputs
    \end{itemize}

    \item \textbf{Deterministic test inputs:}
    \begin{itemize}
        \item Structured plaintext patterns
        \item Fixed keys
    \end{itemize}

    \item \textbf{Cross-validation:}
    \begin{itemize}
        \item Single-threaded, OpenMP, and CUDA outputs were compared
    \end{itemize}
\end{itemize}

No discrepancies were observed, confirming correctness of all implementations.

\section{Performance Evaluation}

\subsection{Experimental Setup}

\begin{itemize}
    \item CPU: AMD Ryzen 7 5800H
    \item GPU: NVIDIA RTX 3050 Ti (95W TDP)
    \item CUDA Toolkit: 12.3
    \item Development Environment: Visual Studio 2022 (MSVC v143)
\end{itemize}

\subsection{Parameters}

Experiments were conducted over varying numbers of plaintext blocks:

\begin{equation}
N = 2^{10} \text{ to } 10^8
\end{equation}

\subsection{Metrics}

\begin{itemize}
    \item Total execution time
    \item Kernel execution time (GPU)
    \item Memory transfer time
    \item Throughput (GB/s)
    \item Speedup over CPU baseline
\end{itemize}

\subsection{Observations}

\begin{itemize}
    \item GPU implementations outperform CPU for large input sizes
    \item Memory transfer dominates performance for small inputs
    \item Bit-slicing improves efficiency for TWINE and WARP
    \item OpenMP provides significant CPU speedup but remains below GPU performance
\end{itemize}

\subsection{Speedup}

Speedup was calculated as:

\begin{equation}
\text{Speedup} = \frac{T_{\text{CPU}}}{T_{\text{GPU}}}
\end{equation}

This was evaluated across all input sizes and cipher implementations.